\makeabstract
{
Angiotensin II Enhances Activity and Flexibility in Murine $CD8^+$ T Cell Immunometabolism
}
{
Audrey Bai (1,2,4), Quinn Hulsey (1), Raghad AlMotairy (1), Allison Norlander (1,2,3)
}
{
\textsuperscript{1}Department of Anatomy, Cell Biology, and Physiology, IU School of Medicine, Indianapolis, IN, USA\\
\textsuperscript{2}Melvin and Bren Simon Comprehensive Cancer Center, IU School of Medicine, Indianapolis, IN, USA\\
\textsuperscript{3}Krannert Cardiovascular Research Center, IU School of Medicine, Indianapolis, IN, USA\\
\textsuperscript{4}Department of Biology, Amherst College, Amherst, MA, USA
}
{
Cytotoxic ($CD8^+$) T cells play a crucial role in adaptive immunity against intracellular pathogens and tumors by killing infected and mutated cells. Persistent antigen exposure, however, can induce an “exhausted” $CD8^+$ phenotype characterized by inhibitory receptors, impaired cytokine production, and poor proliferative potential and survival. As such, exhaustion hinders an effective immune response against cancerous cells, leading to worse disease prognoses and less effective immunotherapies. In this study, we examined $CD8^+$ exhaustion in response to elevated angiotensin II (Ang II) levels, a phenotype implicated in many malignancies. We expected the CD8s of mice with elevated levels of Ang II to exhibit increased exhaustion with a less active and robust metabolic profile, due to the bidirectional relationship between hypertension and chronic inflammation.
}
{
We implanted subcutaneous osmotic pumps containing Ang II or vehicle into 21-day-old male C57BL/6 wild-type mice, incubating for 28 days. We then ran SCENITH analysis of CD8s isolated from harvested mouse spleens. SCENITH (Single Cell ENergetIc metabolism by profilIng Translation inHibition) incorporates flow cytometry to reveal various cell populations’ metabolic profiles.
}
{
The mice of the Ang II treatment group (which were shown to be hypertensive from characteristic fibrosis evident in renal and aortic histology slides) had bulk CD8s that, despite elevated levels of exhaustion markers, exhibited a decreased mitochondrial dependence (MD), a increased glycolytic capacity (GC), and an increased fatty acid oxidation/amino acid oxidation capacity (FAO/AAO cap). When we compared just the exhausted CD8s of the two treatment groups, once again, the hypertensive CD8s exhibited a decreased glycolytic dependence (GD), a decreased MD, an increased GC, and an increased FAO/AAO cap. This was the trend for both examined exhausted CD8 populations.
}
{
Even though Ang II is associated with increased CD8 exhaustion, our results demonstrate its role in enhancing metabolic activity and flexibility in $CD8^+$ T cells, especially in the exhausted subsets. This reveals the need for further elucidation and unraveling of the various effects of Ang II, including its role exacerbating exhaustion, acting as a co-stimulatory molecule in CD8 activation, promoting tumor growth and metastasis, and as a mediator of hypertension and inflammation, etc.
}

\newpage
